\textbf{Principaux domaines abordés :} géométrie dans l'espace

\medskip

\begin{minipage}{0.5\linewidth}
Une maison est modélisée par un parallélépipède rectangle ABCDEFGH surmonté d'une pyramide EFGHS.

On a DC $= 6$,\:\: DA = DH $= 4$. 

Soit les points I, J et K tels que

$\vect{\text{DI}} = \dfrac16\vect{\text{DC}}, \quad \vect{\text{DJ}} = \dfrac14\vect{\text{DA}},\quad \vect{\text{DK}} = \dfrac14\vect{\text{DH}}$.

On note $\vect{\imath} = \vect{\text{DI}},\:\vect{\jmath} = \vect{\text{DJ}},\: \vect{k} = \vect{\text{DK}}$.

On se place dans le repère orthonormé $\left(\text{D}~;~\vect{\imath},\:\vect{\jmath}, \: \vect{k}\right)$.

On admet que le point S a pour coordonnées (3~;~2~;~6).
\end{minipage} \hfill
\begin{minipage}{0.46\linewidth}
\begin{center}
\psset{unit=1cm,arrowsize=2pt 3}
\begin{pspicture}(6.3,6)
\pspolygon(0.3,3.9)(0.3,0.7)(2.7,0.3)(2.7,3.5)
\psline(2.7,0.3)(6,1.1)(6,4.3)(2.7,3.5)
\psline[linestyle=dotted,linewidth=1.25pt](0.3,0.7)(3.6,1.5)(6,1.1)
\psline[linestyle=dotted,linewidth=1.25pt](3.6,1.5)(3.6,4.7)(3.1,5.6)
\psline[linestyle=dotted,linewidth=1.25pt](0.3,3.9)(3.6,4.7)(6,4.3)
\psline(6,4.3)(3.1,5.6)(2.7,3.5)
\psline(3.1,5.6)(0.3,3.9)
\psline(4.3,4.8)(4.3,5.35)
\psline[linewidth=1.25pt]{->}(3.6,1.5)(3.05,1.37)
\psline[linewidth=1.25pt]{->}(3.6,1.5)(4.2,1.4)
\psline[linewidth=1.25pt]{->}(3.6,1.5)(3.6,2.3)
\uput[r](6,1.1){A} \uput[d](2.7,0.3){B} \uput[dl](0.3,0.7){C} \uput[d](3.6,1.5){D}
\uput[ur](6,4.3){E} \uput[dl](2.7,3.5){F} \uput[l](0.3,3.9){G} \uput[ur](3.6,4.7){H}
\uput[u](3.1,5.6){S} \uput[d](4.3,4.8){P} \uput[u](4.3,5.35){Q} \uput[u](3.05,1.37){I}
\uput[ur](4.2,0.9){J} \uput[r](3.6,2.3){K} \uput[d](3.32,1.42){$\vect{\imath}$} \uput[u](3.9,1.3){$\vect{\jmath}$}
\uput[l](3.6,1.9){$\vect{k}$}
\end{pspicture}
\end{center}
\end{minipage}

\medskip

\begin{enumerate}
\item Donner, sans justifier, les coordonnées des points B, E, F et G.
\item Démontrer que le volume de la pyramide EFGHS représente le septième du volume total de la maison.

On rappelle que le volume $V$ d'un tétraèdre est donné par la formule :

\[V = \dfrac13 \times (\text{aire de la base}) \times  \text{hauteur}.\]
\item 
	\begin{enumerate}
		\item Démontrer que le vecteur $\vect{n}$ de coordonnées $\begin{pmatrix}0\\1\\1\end{pmatrix}$ est normal au plan (EFS).
		\item En déduire qu'une équation cartésienne du plan (EFS) est $y + z - 8 = 0$.
	\end{enumerate}	
\item On installe une antenne sur le toit, représentée par le segment [PQ]. On dispose des données suivantes:

\begin{itemize}
\item[$\bullet~~$] le point P appartient au plan (EFS) ;
\item[$\bullet~~$] le point Q a pour coordonnées (2~;~3~;~5,5) ;
\item[$\bullet~~$] la droite (PQ) est dirigée par le vecteur $\vect{k}$.
\end{itemize}
	\begin{enumerate}
		\item Justifier qu'une représentation paramétrique de la droite (PQ) est :

\[\left\{\begin{array}{l c l}
x &=&2\\
y &=&3\\
z &=& 5,5 + t
\end{array}\right. \quad (t \in \R)\]
		\item En déduire les coordonnées du point P.
		\item En déduire la longueur PQ de l'antenne.
	\end{enumerate}	
\item Un oiseau vole en suivant une trajectoire modélisée par la droite $\Delta$ dont une représentation paramétrique est :

\[\left\{\begin{array}{l c r}
x & =& -4 +6s\\ y &=& 7 - 4s\\z &=& 2 + 4s
\end{array}\right.\quad (s \in \R)\]

Déterminer la position relative des droites (PQ) et $\Delta$.

L'oiseau va-t-il percuter l'antenne représentée par le segment [PQ] ?

\end{enumerate}

\bigskip
